\chapter{Publishing a private project}
\label{ch:going-public}

\section{The exercise}

The application described in the previous chapters runs privately, with real
financial data. This repository is a second one, public, built from it.

The goal was not to publish a demo. It was to publish the same application,
adapted so that anybody can run their own instance, with nothing personal left
anywhere: not in the code, not in the documentation, and not in the commit
history.

This chapter describes how that was done, because the process is reusable and
most of its traps are not obvious.

\section{Step one: decide what is personal}

Personal data hides in more places than the obvious ones. The audit found five
categories.

\begin{table}[h]
\centering
\begin{tabular}{@{}p{4.2cm}p{9.5cm}@{}}
\toprule
\textbf{Where} & \textbf{What was found} \\
\midrule
Configuration & The owner's full name as the default value of a report setting. \\
Internal documents & Two reports containing the live application URL, the private
repository link and real transaction counts. \\
Design notes & An internal specification, signed and dated. \\
Scripts & An importer that read the author's old desktop project from a
hardcoded path. \\
Tests & A test that read the author's real broker statements, and skipped when
they were absent. \\
Parsers & A special case with the ISIN of a product the author owns, plus the
names of real people used as text filters. \\
\bottomrule
\end{tabular}
\caption{What the audit found before publishing.}
\end{table}

The first four were removed or replaced. The fifth was rewritten as a generic
pattern, which made the parser both shorter and more capable, as chapter 5
described.

\begin{keypoint}[The useful question]
For each file, ask: \emph{would this still make sense on somebody else's
machine?} A path, a name, an ISIN, a live URL and a count of real records all
fail that question. It catches more than searching for a name.
\end{keypoint}

\section{Step two: a fresh history}

The public repository was started with \code{git init}, not by copying the
private one.

The reason is that git history is data. A repository that was private for months
carries earlier states of every file, including whatever was committed before a
name was moved into an environment variable. Removing a secret from the current
files does nothing if it is still one \code{git log -p} away.

Starting fresh also allowed the history to be rebuilt as a readable narrative:
thirteen commits, each one a feature that works, in the order somebody would
actually build them.

\begin{lstlisting}[style=sh]
chore: bootstrap FastAPI project skeleton
feat: login with bcrypt, signed session and CSRF
feat: expenses and income with monthly and yearly views
feat: broker statement parsers and P&L in Decimal
feat: bank accounts and net worth snapshots
feat: investment positions, price refresh and ISIN lookup
feat: dashboard with net worth, cashflow and category charts
feat: Excel and PDF reports
feat: Telegram bot to log expenses and pull reports
feat: scheduled tasks, health checks and backups over Telegram
chore: initial Alembic migration
build: Docker image, compose stack and deploy blueprints
docs: README, architecture and deployment guides
\end{lstlisting}

Building the history that way exposed a dependency problem. The application
object in \file{app/main.py} grows with every feature, so committing the final
version early would reference routers that did not exist yet, and committing it
late would leave several commits unable to start.

The solution was to reconstruct the intermediate versions: each commit contains
the \file{main.py} that matches the code around it. It costs a few minutes and
produces a history where every commit is a working application.

\section{Step three: the translation}

The private application was written in Spanish: comments, docstrings, variable
names, database columns, routes and interface text. A public repository in one
language and one alphabet is easier for everyone, so everything was translated to
English.

The scale is worth knowing before starting: about seven thousand lines of Python
and eight hundred of templates.

\begin{table}[h]
\centering
\begin{tabular}{@{}p{6.4cm}p{7.3cm}@{}}
\toprule
\textbf{Before} & \textbf{After} \\
\midrule
\code{/gastos}, \code{/inversiones}, \code{/bancos} & \code{/expenses},
\code{/investments}, \code{/banks} \\
\code{importe}, \code{categoria}, \code{tipo} & \code{amount},
\code{category}, \code{kind} \\
\code{valor\_actual}, \code{beneficio}, \code{saldo} & \code{current\_value},
\code{profit}, \code{balance} \\
\code{patrimonio\_snapshots} & \code{net\_worth\_snapshots} \\
\code{"gasto"}, \code{"ingreso"} (stored values) & \code{"expense"},
\code{"income"} \\
\code{/resumen}, \code{/ahorro} (bot) & \code{/summary}, \code{/savings} \\
\bottomrule
\end{tabular}
\caption{Part of the renaming map.}
\end{table}

Two lessons came out of it.

\textbf{Write the map first.} Deciding every rename up front, in a table, turns
the work into mechanical application. Deciding name by name while editing
produces \code{current\_value} in one file and \code{value\_now} in another.

\textbf{Rename the stored values too, not just the columns.} The kind column
held the strings \emph{gasto} and \emph{ingreso}. Renaming the column and leaving
the values would have produced a database that reads half in each language, and
every comparison in the code would have kept a Spanish string literal in it.

\section{Step four: collapse the migrations}

The private repository had three Alembic revisions, describing a history that
included a table for a shared property that no longer exists.

A public repository has no deployment to upgrade, so those three were replaced
by a single initial revision generated from the translated models. The check
that it is correct is two commands:

\begin{lstlisting}[style=sh]
alembic upgrade head     # builds the schema from scratch
alembic check            # no difference between models and database
\end{lstlisting}

\code{alembic check} answering "no new upgrade operations detected" is the proof
that the migration and the models agree.

\section{Step five: rewrite the documentation}

The private repository had a stale README describing two features that had been
removed months earlier, plus two long internal reports full of live URLs.

None of it was reused. The public repository got a README written against the
application as it is today, and two generic guides: an architecture document and
a deployment guide. The internal reports were not adapted; they were replaced,
because a document written for one person and one server does not become general
by deleting the addresses.

\section{Step six: verify, do not assume}

Three checks ran before the repository was considered done.

\textbf{The whole suite.} 118 tests, all passing, after the renaming. A
translation this wide is exactly the kind of change that breaks something
quietly.

\textbf{The application, actually running.} Not just imports: boot the server,
log in, add an expense, see it on the page, download a PDF and a spreadsheet,
and check that every page answers 200.

\begin{lstlisting}[style=sh]
curl -s localhost:8011/health                    # {"status":"ok",...}
curl -s -o /dev/null -w '%{http_code}' localhost:8011/       # 303 -> /login
# after logging in: /expenses, /investments, /banks, /reports  -> 200
# report download                                              -> 200 application/pdf
\end{lstlisting}

\textbf{A final search for anything personal.} One \code{grep} over the whole
repository for names, old route words, old domain names and old table names. The
only remaining hits should be the ones that are supposed to be there: the
repository URL and the copyright line of the licence.

\begin{pitfall}[The one that nearly went wrong]
The default SSH key on the machine authenticated as a different GitHub account,
which had no access to the target one. The first push would have failed with a
confusing "repository not found". Checking \code{ssh -T} and configuring a host
alias took two minutes; debugging it later, with a repository half pushed, would
have taken longer.
\end{pitfall}

\section{A checklist}

\begin{enumerate}
  \item Audit for anything that only makes sense on your machine.
  \item Start a new repository; never publish a history that was private.
  \item Write the renaming map before renaming anything.
  \item Collapse migrations into one initial revision.
  \item Rewrite the documentation for a stranger, do not adapt internal notes.
  \item Rebuild the history as commits that each work.
  \item Run the tests, run the application, then grep for what should not be
        there.
\end{enumerate}
